\documentclass[10pt,a4paper]{amsart}
\usepackage[margin=19mm]{geometry}
\usepackage{amsmath,amssymb,mathtools}
\usepackage[scr=rsfs,cal=euler]{mathalfa}
\usepackage{xfrac}
\usepackage[unicode,hidelinks,bookmarks=false]{hyperref}
\newcommand{\ie}{i.e.\ }
\newcommand{\cf}{cf.\ }
\newcommand{\resp}{respectively\ }
\newcommand{\Subsection}{\subsection}
\providecommand{\nobreakdash}{\nobreak\hbox{-}}
\setlength{\emergencystretch}{2em}
\multlinegap0pt
\usepackage{amssymb}
\usepackage{dsfont}
\usepackage{tikz}
\usepackage{mathtools}
\usepackage{bm}
\usepackage[shortlabels]{enumitem}
\usepackage[normalem]{ulem}
\newtheorem{proposition}{Proposition}
\newcommand{\R}{{\mathbb R}}
\title{Hydrodynamics and boundary-induced phase transitions in the $n$-species particle-exchange process}
\author{Gunter M. Schutz, Ali Zahra}
\date{Source: arXiv:2605.07615v1 (CC BY 4.0)}
\begin{document}
\maketitle

\noindent\emph{Source and licence.} Hydrodynamics and boundary-induced phase transitions in the $n$-species particle-exchange process. Version arXiv:2605.07615v1 (CC BY 4.0), distributed by arXiv under the Creative Commons Attribution 4.0 International licence, http://creativecommons.org/licenses/by/4.0/, as recorded in the arXiv licence metadata for this exact version and shown on the abstract page https://arxiv.org/abs/2605.07615v1. Full licence notice: \texttt{licenses/arxiv-2605.07615-CC-BY-4.0.txt}.\par\smallskip

\noindent\textit{Editorial scope.} Counted unit: the article's proposition proposition (source0.tex lines 705--711). The article's complete original proof of it is delivered with it (source0.tex lines 713--729, one \texttt{proof} environment of 17 lines). No mathematics is written, repaired or re-derived: the statement and the proof are the article's own lines, in the article's own notation, and no retained line is substituted or re-flowed.  Retained uncounted as the source-local context the proof needs: lines 695--698.  Nothing was omitted from any retained span, so the package carries no omission record.  No retained line contains a citation, so the wrapper carries no bibliography and no bibliography entry.  No retained line contains a figure environment, an image-inclusion command or drawing code, so no image is delivered and the package declares no asset dependency.  Editorial formatting only, in addition: an AMS \texttt{amsart} preamble that restates the article's own declarations and macros used by the retained lines, namely the environments \texttt{proposition} on separate counters, together with the standard packages those lines need.  The retained environments print in this document as Proposition 1 in order of appearance, whereas the article prints the same items with the numbering of its own full document.  The complete native source of the article as distributed in the arXiv e-print for this exact version is bundled unchanged as \texttt{sources/200095/source0.tex}.
\begin{equation}
R_n(\omega) := \sum_{\alpha=0}^n \frac{\rho_\alpha}{f_\alpha-\omega}.
\label{eq:Rmu}
\end{equation}

\begin{proposition}[Existence of one root in each gap]
For every state $\rho\in\mathring{\mathcal D}$, the function $R_n$ has exactly one simple zero in each interval
\[
\bigl(f_{\sigma(i-1)},f_{\sigma(i)}\bigr),
\qquad i=1,\dots,n.
\]
\end{proposition}

\begin{proof}
Differentiate \eqref{eq:Rmu} with respect to $\omega$:
\[
R_n'(\omega)=\sum_{\alpha=0}^n \frac{\rho_\alpha}{(f_\alpha-\omega)^2}>0
\]
whenever $\omega\neq f_\alpha$ for all $\alpha$. Hence $R_n$ is strictly increasing on each connected component of
\[
\R\setminus\{f_0,\dots,f_n\}.
\]
Now fix $i\in\{1,\dots,n\}$. As $\omega\downarrow f_{\sigma(i-1)}$, the term with index $\sigma(i-1)$ dominates and
$R_n(\omega)\to -\infty.$
As $\omega\uparrow f_{\sigma(i)}$, the term with index $\sigma(i)$ dominates and
$R_n(\omega)\to +\infty.$
Since $R_n$ is continuous and strictly increasing on the interval
$\bigl(f_{\sigma(i-1)},f_{\sigma(i)}\bigr),$
it has exactly one zero there. The non-degeneracy of the root follows since $R_n'(\omega)>0$ at the zero.
\end{proof}

\end{document}
